\section{Making review binding}
\label{sec:14.3}
A finding binds when it changes the institution's ability to continue without requiring the institution to agree. The working model is the federalwide assurance in human-subjects research: the board doesn't punish, it withdraws clearance, and a separate federal office can suspend the assurance, putting funding eligibility at risk. Each link matters — jurisdiction before the disputed protocol, access the operator can't withhold, authority to suspend, and a licensing body \emph{obliged} to act rather than invited to reconsider. If another discretionary decision by the reviewed institution is required anywhere in the chain, the ruling is advisory at that point.

The body that does this for the floor is what the rest of the book calls the floor court: an Article III court, meaning a federal court whose judges hold office under Article III of the United States Constitution, during good behavior and at a salary that can't be cut, whose jurisdiction is fixed before the deployment it hears, with access to the record the operator can't withhold, power to grant interim relief, and remedies, a finding of wrongful refusal and the forfeiture of a bond, that take effect without the operator's agreement. Appendix~\ref{sec:B} says how it would be constituted.
